\documentclass[11pt]{article}
\usepackage[a4paper, top=18mm, bottom=18mm, left=20mm, right=20mm]{geometry}
\usepackage[T2A]{fontenc}
\usepackage[utf8]{inputenc}
\usepackage[ukrainian]{babel}
\usepackage{graphicx}
\usepackage[export]{adjustbox}
\usepackage{amsmath}
\usepackage{amssymb}
\graphicspath{{/app/Quiz/Scripts/2023/ProgAll/15BinTree/}{/app/Quiz/Scripts/2021/Mod2021_3/BinTree/}{/app/Quiz/Scripts/2020/Mod2020_4/BinTree/}}
\setlength{\parindent}{0pt}
\setlength{\parskip}{6pt}
\pagestyle{empty}
\begin{document}
%begin
{\large\textbf{Контрольна робота}}

\textbf{Варіант \No\,10}\hfill \textit{ПІБ:}\,\underline{\hspace{60mm}}
\vspace{4mm}

%beginitem
1. Що буде виведено за виконання фрагменту коду?

%code
#include <iostream>

int h(int a, int b){
    int c = 58;
    if (a < 2)
        return 6;
    if (a > 5)
         c = 2;
    else 
        return 7;
    return c;
}

int main(){
    std::cout << h(-2, -2);
    return 0;
}
%code
\vspace{5mm}
%enditem

%beginitem
2. Що буде виведено за виконання фрагменту коду?
try:
    res = 8<True
    if isinstance(res, bool):
        print(f'bool;{res}')
    elif isinstance(res, int):
        print(f'int;{res}')
    elif isinstance(res, float):
        print(f'float;{res:.2f}')
    else:
        print('???')
except:
    print('Z')
\vspace{5mm}
%enditem

%beginitem
3. %goodpath:Scripts/2018/Mod2018_1/03-good.txt
Відмітити все, що є коректним оператором С++:
( 1 ) bool f(int a);

( 2 ) int f(int a,b);

( 3 ) cout<<3.4;

( 4 ) int n; bool n;
\vspace{5mm}
%enditem

%beginitem
4. 

\begin{center}
	\adjustimage{max size={0.8\linewidth}{0.8\paperheight}}
	{../BinTree/trees_exp/178.png}
\end{center}
Указати послідовність міток вершин дерева, зображеного на рис., що утвориться в результаті обходу дерева в прямому порядку. Мітки записувати через пробіл.\par Обхід дерева починається з кореня (на рис. це верхня вершина).
%answer:178 прямому
\vspace{5mm}
%enditem

%beginitem
5. Що буде виведено за виконання фрагменту коду?
def f(n):
    print('f', end='')
    return n

try:
    print(f(-2) and f(1))
except:
    print('ERR')
\vspace{5mm}
%enditem

%beginitem
6. Що буде виведено за виконання фрагменту коду?
g=[10,11,12,13,14]; del g[6:]; print(g)\vspace{5mm}
%enditem

%beginitem
7. Що буде виведено за виконання фрагменту коду?

%code
print(5/5*5*5)
%code
\vspace{5mm}
%enditem

%beginitem
8. Що буде виведено за виконання фрагменту коду?
try:
    print(6, end=';')
    print(2j<6j, end=';')
    print(7, end=';')
except TypeError:
    print(9, end=';')
except ValueError:
    print(4, end=';')
except:
    print('Z', end=';')
else:
    print(3, end=';')
finally:
    print(0)
\vspace{5mm}
%enditem

%beginitem
9. %goodpath:Scripts/2019/Mod2019_1/Lex/03-good.txt
Відмітити все, що є коректним оператором С++:
( 1 ) 3.14%2

( 2 ) if a<b: a=0 else: b=0

( 3 ) i,j=j,j

( 4 ) x+3=5
\vspace{5mm}
%enditem

%beginitem
10. %goodpath:Scripts/2018/Mod2018_1/01-good.txt
Відмітити все, що є однією правильною лексемою:
( 1 ) Switch

( 2 ) 0,5

( 3 ) 8.3

( 4 ) 8,3
\vspace{5mm}
%enditem

%end
\newpage
\end{document}

